\subsection{The exact lifted link of the selected-cusp normalization}
\begin{statement}{P3F.55}{Proved}
Let
\[
 f(s,\rho)=(X,Y,R)
       =(s^2-\rho^2,\ s(s^2-\rho^2),\ \rho):
                    \mathbb C^2\longrightarrow\mathbb C^3.       \tag{CL.1}
\]
Give both complex spaces their complex orientations. Choose a small positive
\(\epsilon\), put \(L_\epsilon=f^{-1}(S^5_\epsilon)\), and give this
three-manifold the outward boundary orientation of
\(f^{-1}(B^6_\epsilon)\). Choose a point of \(S^5_\epsilon\) outside its
image and an orientation-preserving stereographic chart
\(\psi:S^5_\epsilon\setminus\{p\}\to\mathbb R^5\).
For any positive constant \(\eta\), define
\[
 h(s,\rho)=\eta\,
       \frac{\operatorname{Re}(s\overline\rho)}
            {\sqrt{|s|^2+|\rho|^2}},\qquad
 e=(\psi f,h):L_\epsilon\longrightarrow\mathbb R^5\times\mathbb R.
                                                                  \tag{CL.2}
\]
Then \(L_\epsilon\) is an oriented three-sphere, \(e\) is a smooth
embedding, and its Haefliger invariant in the convention of
Gauniyal--Turchin, Theorem 4.1, is \(+1\).

This statement computes the specified local lifted link. Its comparison with a completed source-parameter family requires the separate ambient and circle-line morphisms.
\end{statement}
\begin{proof}
\par

\medskip\noindent\textbf{1. The sphere, the immersion, and its complete double locus}

\par\nobreak

The function
\[
 \|f(s,\rho)\|^2
   =|s^2-\rho^2|^2+|s(s^2-\rho^2)|^2+|\rho|^2
\]
is proper and has strictly positive derivative along every positive radial
ray away from the origin: its three summands scale by \(r^4,r^6,r^2\),
respectively, and cannot vanish simultaneously away from zero. Each ray
therefore meets \(L_\epsilon\) once and transversely. Radial projection
is a smooth orientation-preserving diffeomorphism
\(L_\epsilon\to S^3\), with the boundary orientations just specified.

The actual derivative is
\[
 df=\begin{pmatrix}
       2s&-2\rho\\
       3s^2-\rho^2&-2s\rho\\
       0&1
     \end{pmatrix}.                                               \tag{CL.3}
\]
Its minor ideal is \((s,\rho^2)\); hence its only rank drop is at the
origin. The radial transversality above implies that
\(f|_{L_\epsilon}\) is an immersion into \(S^5_\epsilon\).

Two points with the same image have the same \(\rho\) and the same
\(s^2\). The unequal pair is \(s,-s\), and the equality of the second
coordinates is \(2s(s^2-\rho^2)=0\). Thus the complete double locus is
\[
 L_+=\{s=\rho,\ |\rho|=\epsilon\},\qquad
 L_-=\{s=-\rho,\ |\rho|=\epsilon\}.                                \tag{CL.4}
\]
Its image is \(D=\{(0,0,\rho):|\rho|=\epsilon\}\).
There are no triple points. At a double point the two transverse complex
directions are \((2\rho,2\rho^2,0)\) and
\((-2\rho,2\rho^2,0)\). They are independent for \(\rho\ne0\);
the sheets meet transversely along their common complex \(\rho\)-line.
Their restrictions to the links therefore meet transversely along \(D\).
The projected immersion is generic in the precise sense of the theorem.

On the two double branches \(h=\pm\eta|\rho|/\sqrt2\).
Consequently (CL.2) separates every pair not separated by \(f\).
It is an injective immersion of a compact manifold, hence an embedding.
This proves the actual lift, including its ordering.

\par

\medskip\noindent\textbf{2. The double-link orientation and the exact pushed circle}

\par\nobreak

We use the outward-normal-first boundary convention and the transverse
intersection orientation specified in Section 1 of Gauniyal--Turchin.
For the two complex sheets in \(\mathbb C^3\), their common complex
line has its complex orientation: the complementary complex directions
give a positive real determinant. On the link the outward common radial
vector precedes its phase vector. Thus the Ekholm orientation of \(D\)
is increasing \(\arg\rho\), and the induced orientations on both
\(L_+\) and \(L_-\) are also increasing \(\arg\rho\).
The two transverse complex lines \(s=\rho\) and \(s=-\rho\) meet at
the origin with intersection number \(+1\). Equivalently their radial
images form the positively oriented Hopf link. Hence
\[
                         \operatorname{lk}(L_+,L_-)=1.            \tag{CL.5}
\]

The theorem's invariant \(E\) uses a normal field on the entire double
preimage whose pushed union has zero class in the complement. We give
that field explicitly. On either branch set
\[
                         v_\pm=\overline\rho\,\partial_s.         \tag{CL.6}
\]
These vectors are tangent to \(L_\epsilon\): at \(X=Y=0\), the
differential of \(\|f\|^2\) depends only on the real radial
\(\rho\)-coordinate, which (CL.6) leaves fixed. They are transverse
to the double circle. The constant \(\partial_s\) framing has
self-linking zero. Indeed the holomorphic coordinates
\((s\mp\rho,\rho)\) identify its small push with a parallel complex
disc disjoint from the core disc. Multiplication by \(\overline\rho\)
has winding \(-1\), so the self-linking of each pushed branch is
\(-1\). Together with (CL.5), each total linking is
\(-1+1=0\), exactly the zero-class requirement of their equations
(4.1)--(4.2).

No normal field is now replaced by a homotopic one without its derivative.
The actual two images of (CL.6) satisfy
\[
 df_{\rho,\rho}(v_+)+df_{-\rho,\rho}(v_-)
                    =(0,4|\rho|^2\rho,0).                       \tag{CL.7}
\]
The pushed image circle is therefore, after its positive radial
readjustment, isotopic in the complement to the link of the complex line
\[
                         (X,Y,R)=(0,c\rho,\rho),\qquad c>0.       \tag{CL.8}
\]
The image of (CL.1) is the reduced hypersurface
\[
                         Y^2-X^2(X+R^2)=0.                      \tag{CL.9}
\]
Its normalization is exactly (CL.1): \(s\) is integral, satisfying
\(s^2=X+R^2\), and \(s=Y/X\) on the dense set \(X\ne0\).
On the disc bounded by (CL.8), the left side of (CL.9) is
\(c^2\rho^2\). Its winding is \(+2\). A small transverse analytic
perturbation gives the same two positive intersections, so the oriented
linking number of the pushed circle and the immersed link is
\[
                                  E=2.                           \tag{CL.10}
\]
This computes the specific tangent push in the cited formula; it does
not substitute a differently framed normal push.

\par

\medskip\noindent\textbf{3. The actual Stiefel and quaternion comparison}

\par\nobreak

Némethi--Pintér, Theorem 1.2.2, Proposition 4.2.3 and Proposition 6.1.1,
apply to the holomorphic germ (CL.1), whose isolated rank drop was proved
above. The quotient of its minor ideal has basis \(1,\rho\). Their
complex Smale class is therefore \(2[u]\), and its realification is
\[
                  \Omega(f)=-2[L]\in\pi_3(SO(5)).                \tag{CL.11}
\]
Here we retain their literal convention
\(q=z+w j\), with \((z,w)\) carrying the complex orientation,
\(L_q(p)=qp\), and \(u_q(p)=pq\) in its complex two-dimensional
matrix representation. Their realification arrow sends \([u]\) to
\(-[L]\). The local stereographic chart is the same kind of
oriented link chart used in Proposition 6.1.1, so (CL.11) concerns the
actual projected link, with no independent source framing substituted.

We next compute the morphism needed by the coefficient
\(\Omega_{S^2}\) in Gauniyal--Turchin. The fibre inclusion
\(S^2\to V_{5,3}\) is the orbit of the last vector under a fixed
three-plane rotation. Its quaternion lift is
\(\operatorname{Ad}_q(p)=qpq^{-1}\). Proposition 4.2.2(c) of
Némethi--Pintér gives \([\operatorname{Ad}]=2[L]\) stably.
The projection of this lift is
\[
                            q\longmapsto qi q^{-1}.              \tag{CL.12}
\]
Its Hopf invariant is \(-1\) in the specified complex source
orientation. To check the sign directly, its circle fibres are
\[
          (z,w)\longmapsto(ze^{i\theta},we^{-i\theta}).
\]
At \(q=1\) the derivative on the complementary \(w\)-plane is
\((u,v)\mapsto(2v,-2u)\) in the ordered \(j,k\)-plane, with
positive determinant. Thus the oriented fibre over \(i\) is the
positive \(z\)-circle. The oriented fibre over \(-i\) is the
negative \(w\)-circle, by the same right-action vector \(qi\).
Their linking number in the complex-oriented three-sphere is \(-1\),
which proves the assertion.

The Whitehead square \([\operatorname{id}_{S^2},
\operatorname{id}_{S^2}]\) has Hopf invariant \(+2\): the attaching
map of the oriented top cell of \(S^2\times S^2\), followed by
the fold of its two two-spheres, has cup square
\((a+b)^2=2ab\). The Hopf fibration exact sequence identifies
\(\pi_3(S^2)\) with \(\mathbb Z\) by this linking invariant.
Consequently this precise Whitehead square is \(-2\) times
(CL.12), and its image in the Stiefel group is
\[
       [\operatorname{id},\operatorname{id}]
                \longmapsto-2[\operatorname{Ad}]=-4[L].         \tag{CL.13}
\]
Combining (CL.11) and (CL.13), with both arrows typed, gives
\[
                             \Omega_{S^2}(f)=\frac12.             \tag{CL.14}
\]

\par

\medskip\noindent\textbf{4. The embedding invariant and its exact scope}

\par\nobreak

All hypotheses of Gauniyal--Turchin, Theorem 4.1, have been verified:
the lift is an actual embedding; its projection is a generic immersion;
both double-preimage components are explicitly oriented; their required
total-link-zero framing and its push have been computed; and the
Stiefel coefficient uses their Whitehead generator. The spherical
version stated directly before that theorem applies after the chosen
stereographic chart. Its formula gives
\[
 \mathcal H(e)
   =\frac12\operatorname{lk}(L_+,L_-)
      +\frac16E(f)+\frac13\Omega_{S^2}(f)
   =\frac12+\frac13+\frac16=1.                                  \tag{CL.15}
\]
Changing \(\eta\) through positive values is a literal isotopy of
embeddings, since the projected double pairs remain separated.
The same holds for any interpolation of the height that preserves
the displayed two-branch order.

This proves the invariant of (CL.2). The ambient source-to-orbit map,
the outside cap, the circle Euler line, and the boundary-normal
comparison of a completed threefold are additional morphisms.
Their proof belongs to P3F.56; (CL.15) alone does not identify that
threefold's return.

\par

\medskip\noindent\textbf{5. Exact comparison with the corrected Seifert-surface formula}

\par\nobreak

We identify the normalization of \(\mathcal H\) just computed with
the corrected Takase invariant, rather than infer a sign from their
common name. Both use the outward-normal-first boundary convention.
Takase's correction paper, Corollary 2.6, gives the group isomorphism
\[
 \mathcal H_{\rm T}(e)
      =\frac{e(\nu V)^2-\sigma(V)}8                 \tag{CL.16}
\]
for an actual oriented Seifert embedding \(V^4\subset S^6\).
Gauniyal--Turchin, Theorem 4.1, gives the group isomorphism used in
(CL.15). We compare them on an explicitly nonzero embedded projection.

Takase, Remark 4.2, gives a punctured Kummer surface
\(V\subset S^5\), with \(\sigma(V)=-16\); its existence is
also the example of Haefliger (1966), Theorem 5.17. Its normal line
in \(S^5\) is trivial because it is oriented. Composing with the
equatorial inclusion in \(S^6\) therefore gives a trivial normal
two-plane, so \(e(\nu V)=0\). Choose a point of \(S^5\) outside
this compact surface and an orientation-preserving stereographic
chart. The boundary becomes an actual embedding
\(g:S^3\hookrightarrow\mathbb R^5\), with empty double locus,
and the same Seifert surface in \(\mathbb R^5\).

Use the oriented boundary identification of Némethi--Pintér, Section
2.4, with the same outward-normal-first convention. Their sign-refined
Theorem 9.1.6, applied to the embedded Seifert surface, gives
\[
                 \Omega_a(g)=\frac32\sigma(V)=-24
                         \quad\hbox{in the }[L]\hbox{ convention}.
                                                               \tag{CL.17}
\]
The literal fibre-inclusion morphism (CL.13) sends the Whitehead
square to \(-4[L]\). It follows that
\(\Omega_{S^2}(g)=6\). Since the double locus is empty, the two
linking terms in Gauniyal--Turchin's formula vanish, and hence
\[
                  \mathcal H(g)=6/3=2
                    =\frac{0-(-16)}8=\mathcal H_{\rm T}(g).
                                                               \tag{CL.18}
\]
In particular this is a nonzero element of the embedding group.
The two displayed invariants are isomorphisms of that same group
\(C_3^3\cong\mathbb Z\). Their composite automorphism of
\(\mathbb Z\) fixes the nonzero integer two; the only such
automorphism is the identity. Thus
\[
                           \mathcal H=\mathcal H_{\rm T}
                                                               \tag{CL.19}
\]
on the entire group, with its specified orientation convention.
This identifies the exact local value \(+1\) in (CL.15) with
the corrected Seifert-surface formula. It still makes no inference
about a source-family return without the separate morphisms of
P3F.56. As an additional direct source check, Sakai (2015),
Proposition 6.8 and its proof, state precisely the same embedded-
projection comparison \(\mathcal H(g)=-\Omega_a(g)/12
=-\sigma(V)/8\).

The exact certificates are \path{code/part3_root/source_cusp_link.py} and
\path{code/part3_root/haefliger_normalization_bridge.py}, with their JSON and
captured output. The second verifies every integer and rational arrow
in the nonzero normalization test (CL.17)--(CL.19). The local-link
certificate also checks the perturbation
\((s^2-\rho^2,s(s^2-\rho^2+\varepsilon),\rho)\): it has two ordinary
crosscaps at \(s=0,\rho^2=\varepsilon\) and no triple point.
The first structural-factor comparison failure is preserved in
\path{code/part3_root/attempts/}; its correction uses an exact polynomial
zero-difference check.
\end{proof}
\par\medskip\noindent\textbf{Verified primary sources.} \href{https://doi.org/10.4171/CMH/363}{N\'emethi--Pint\'er (2015)}, \S2.4, Propositions 4.2.2--4.2.3 and 6.1.1, and Theorem 9.1.6; \href{https://doi.org/10.1017/S0013091526101485}{Gauniyal--Turchin (2026)}, \S1, \S2.1, equations (4.1)--(4.3), and Theorem 4.1 with its preceding spherical-knot qualification; \href{https://arxiv.org/abs/math/0506356}{Takase, correction preprint (2005)}, Corollary 2.6 and Remark 4.2; \href{https://www.maths.ed.ac.uk/~v1ranick/tda/haefliger.pdf}{Haefliger (1966)}, Theorem 5.17, printed p.421, and its preceding compression proof; \href{https://doi.org/10.4310/HHA.2015.v17.n2.a15}{Sakai (2015)}, Proposition 6.8.
